% Pursuing Stacks (cotes 134-2 to 134-8): archive cote 134-5, site page 35 (PDF page 36), typescript page 222.
% typescript. Draft compared: ps4.tex:1687-1754.
% AI-transcribed by Claude Opus 5.5 (claude-opus-5-5, adaptive thinking, effort high), Message Batches API, batch ps2 (msgbatch_01NVwr1Yuo8qJf5rykeGBXPK), 2026-11-12; not yet checked by eye.
\scanpage{134-5/35}
\mspage{222}
generalities on presheaves on a product category $A\mathrm{x}B$ , where for the
time being $A,B$ are any two objects of (Cat). The following notation
is often useful, for a pair of presheaves
\[ F\in \mathrm{Ob}\, A\hat{} \quad , \qquad G\in \mathrm{Ob}\, B\hat{} \quad , \]
introducing an``\uline{external product}''
\[ F\boxtimes G \in \xout{\illegible{1}}\mathrm{Ob}(A\mathrm{x}B)\hat{} \]
by the formula
\[ F\boxtimes G\ (a,b) = F(a)\mathrm{x}G(b) \ . \]
Introducing the two projections
\[ p_1 : A\mathrm{x}B \longrightarrow A \quad , \qquad p_2 : A\mathrm{x}B \longrightarrow B \quad , \]
and the corresponding inverse image functors $p_i^{*}$ , we have
\[ F\boxtimes G \ = \ p_1^{*}(F)\mathrm{x}p_2^{*}(G) \qquad . \]
Of course $F\boxtimes G$ depends bifunctorially on the pair $(F,G)$, and it is ea-%
sily checked\hand{,}by the way\hand{,}that the corresponding functor
\[ A\hat{}\mathrm{x}\ B\hat{} \longrightarrow (A\mathrm{x}B)\hat{} \quad , \qquad (F,G)\longmapsto F\boxtimes G \quad , \]
is fully faithful.

We have a tautological canonical isomorphism
\[ (A\mathrm{x}B)_{/F\boxtimes G} \xrightarrow{\sim} A_{/F}\ \mathrm{x}\ B_{/G} \qquad , \]
and hence the

\uline{Proposition}\hand{1}\unsure{\hand{:}} a)If $F$ and $G$ are aspheric objects in $A\hat{}$and $B\hat{}$\hand{/}respectively,
then $F\boxtimes G$ is an aspheric object of $(A\mathrm{x}B)\hat{}$.
b) If $u : F\rightarrow F'$ and $v : G\rightarrow G'$ are aspheric maps in $A\hat{}$and $B\hat{}$ res-%
pectively, then
\[ u\boxtimes v \ : \ F\boxtimes G \rightarrow F'\boxtimes G' \]
is an aspheric map in $(A\mathrm{x}B)\hat{}$.

Indeed, this follows respectively from the fact that a product of
two aspheric objects (resp. aspheric maps) in (Cat) is again aspheric.

\uline{Corollary} 1\unsure{\hand{:}}Let $F$ in $A\hat{}$ be aspheric over the final object $e_{A\hat{}}$ , then
$F\boxtimes e_{B\hat{}} = p_1^{*}(F)$ is aspheric over the final object of $(A\mathrm{x}B)\hat{}$.

\inserted[typed]{\xout{\illegible{1}}}\uline{Corollary 2.} Assume $A$ and $B$ are totally aspheric, then so is $A\mathrm{x}B$.

We have to check\xout{s} that the product elements
\[ (a,b)\mathrm{x}(a',b') \ = (a\mathrm{x}a')\boxtimes(b\mathrm{x}b') \]
are aspheric, which follows from the assumption (namely $a\mathrm{x}a'$ and $b\mathrm{x}b'$
aspheric) and part a) of the proposition.

\margin{35}
